% Standalone resolution manuscript generated from solution.md.
% Compile with XeLaTeX; author and review metadata are maintained in Markdown.
\documentclass[10pt,a4paper]{article}
\usepackage[margin=24mm,headheight=15pt]{geometry}
\usepackage{fontspec}
\setmainfont{DejaVuSerif.ttf}[BoldFont=DejaVuSerif-Bold.ttf,ItalicFont=DejaVuSerif-Italic.ttf,BoldItalicFont=DejaVuSerif-BoldItalic.ttf]
\setsansfont{DejaVuSans.ttf}[BoldFont=DejaVuSans-Bold.ttf,ItalicFont=DejaVuSans-Oblique.ttf]
\setmonofont{DejaVuSansMono.ttf}
\usepackage{amsmath,amssymb,unicode-math}
\setmathfont{latinmodern-math.otf}
\usepackage{xcolor,microtype,parskip,needspace,fancyhdr}
\usepackage{longtable,booktabs,array}
\newcounter{none} % Pandoc's unnumbered tables
\usepackage{bookmark}
\usepackage{xurl}
\hypersetup{colorlinks=true,urlcolor=blue,linkcolor=blue,pdftitle={IE-05:
An order-eight counterexample to the orthogonal extremizer
conjecture},pdfauthor={George Stepaniants},pdflang={en-GB}}
\urlstyle{same}
\setlength{\emergencystretch}{3em}
\providecommand{\tightlist}{\setlength{\itemsep}{0pt}\setlength{\parskip}{0pt}}
\providecommand{\pandocbounded}[1]{#1}
\setcounter{secnumdepth}{0}
\allowdisplaybreaks[1]
\widowpenalty=10000
\clubpenalty=10000
\pagestyle{fancy}
\fancyhf{}
\fancyhead[L]{\small\sffamily RESOLUTION / IE-05}
\fancyhead[R]{\small\sffamily 11 September 2026}
\fancyfoot[L]{\parbox[b]{0.9\textwidth}{\fontsize{8}{10}\selectfont AI-assisted
proof; independent automated-agent review documented in the submission
record.}}
\fancyfoot[R]{\footnotesize\thepage}
\begin{document}
{\raggedright\Large\bfseries IE-05: An order-eight counterexample to the
orthogonal extremizer conjecture\par}
\vspace{0.4em}
{\normalsize\bfseries George Stepaniants\par}
{\small Department of Computing and Mathematical Sciences, California
Institute of Technology, Pasadena, California, USA.\par}
\vspace{0.6em}
\pagestyle{plain}
\setlength{\abovedisplayskip}{6pt}
\setlength{\belowdisplayskip}{6pt}
\setlength{\abovedisplayshortskip}{4pt}
\setlength{\belowdisplayshortskip}{4pt}

This manuscript gives a negative resolution of the exact extremizer
equality in IE-05. It was developed with substantial ChatGPT/Codex
assistance at the author's request. A separate
\href{https://github.com/ajt60gaibb/OpenProblemsInNLA/blob/main/references/stepaniants-ie05-2026-09-11/independent-review.md}{independent
Codex-agent review} returned \textbf{PASS for the complete canonical
target}, with exact rational checks of every pivot and active entry.
This is automated-agent review, not external human peer review or formal
proof-assistant verification.

\subsection{Theorem}\label{theorem}

Let \(L_8\) be the unit lower triangular matrix with every entry below
the diagonal equal to \(-1\), and let \(Q_8\) be its orthogonal factor
in the QR factorization with positive diagonal in the triangular factor.
There is an explicit real orthogonal matrix \(\widetilde Q\in O(8)\) for
which partial pivoting, choosing the first available row in every tie,
gives

\nopagebreak[4]

\[
\rho_{\mathrm{PP}}(\widetilde Q)=\frac{5272}{63}
>\sqrt{\frac{17948132}{2601}}
=\rho_{\mathrm{PP}}(Q_8).
\]

In particular, the equality proposed in IE-05 is false. The supremum
over all orthogonal matrices and admissible partial-pivoting paths is
strictly greater than the proposed value already at order eight.

All computations below are in exact arithmetic. The proof supplies
integer matrices, rational squared norms and exact pivot paths. Decimal
approximations are unnecessary for the comparison.

\subsection{1. The matrix and its QR
convention}\label{the-matrix-and-its-qr-convention}

Put

\nopagebreak[4]

\[
\widetilde L=L_8+e_8e_2^T.
\]

Thus the only modification is that the entry in row \(8\), column \(2\)
is changed from \(-1\) to \(0\). Define

\nopagebreak[4]

\[
H=
\begin{pmatrix}
 1&-1&-3&-21&-41&-101&-325&63\\
-1& 3& 1& 17& 71& 291&1179&31\\
-1&-1&13&-19&-56&-196&-752&16\\
-1&-1&-3&189&-28&-98&-376&8\\
-1&-1&-3&-51&581&-49&-188&4\\
-1&-1&-3&-51&-213&1507&-94&2\\
-1&-1&-3&-51&-213&-873&2589&1\\
-1& 1&-5&-55&-183&-683&-2683&1
\end{pmatrix}
\]

and

\nopagebreak[4]

\[
D=\operatorname{diag}(8,16,240,47640,472430,3644970,16148136,5272).
\]

Direct integer multiplication gives

\nopagebreak[4]

\[
H^TH=D.
\]

Consequently

\nopagebreak[4]

\[
\widetilde Q=HD^{-1/2}
\]

is exactly orthogonal. To identify its QR convention and elimination
multipliers, define the upper triangular matrix

\nopagebreak[4]

\[
T=
\begin{pmatrix}
1&-1&-3&-21&-41&-101&-325&63\\
0&2&-2&-4&30&190&854&94\\
0&0&8&-44&-67&-107&-223&173\\
0&0&0&120&-106&-116&-70&338\\
0&0&0&0&397&-183&48&672\\
0&0&0&0&0&1190&190&1342\\
0&0&0&0&0&0&3063&2683\\
0&0&0&0&0&0&0&5272
\end{pmatrix}.
\]

Another direct integer multiplication gives

\nopagebreak[4]

\[
H=\widetilde L T.
\]

Therefore

\nopagebreak[4]

\[
\widetilde L=\widetilde Q\,(D^{1/2}T^{-1}).
\]

The second factor is upper triangular with positive diagonal, so
\(\widetilde Q\) is precisely the orthogonal factor in the specified
positive-diagonal QR factorization of \(\widetilde L\).

\subsection{2. An admissible partial-pivoting
path}\label{an-admissible-partial-pivoting-path}

The factorization

\nopagebreak[4]

\[
\widetilde Q=\widetilde L\,(TD^{-1/2})
\]

is its no-exchange LU factorization, with positive nonzero diagonal in
the upper factor. After eliminating columns \(1,\ldots,k-1\), the active
first-column entries are

\nopagebreak[4]

\[
(\widetilde S_k)_{i,k}
=\widetilde L_{i,k}\,(TD^{-1/2})_{k,k},
\qquad i=k,\ldots,8.
\]

Because \(\widetilde L_{k,k}=1\) and every strict-lower entry of
\(\widetilde L\) is \(-1\) or \(0\), the diagonal entry is a
largest-magnitude entry in the active first column. It is also its first
available row. Thus the stipulated tie rule makes no row exchanges at
any step. The elimination multipliers are exactly the strict-lower
entries of \(\widetilde L\).

More explicitly, if \(I_k=\{k,k+1,\ldots,8\}\), then the active Schur
complement is

\nopagebreak[4]

\[
\widetilde S_k
=\widetilde L_{I_k,I_k}T_{I_k,I_k}D_{I_k,I_k}^{-1/2}.
\]

This formula reduces every squared entry comparison to rational
arithmetic.

\subsection{3. Exact growth of the
counterexample}\label{exact-growth-of-the-counterexample}

The largest absolute entries of the eight columns of \(H\) are,
respectively,

\nopagebreak[4]

\[
(1,3,13,189,581,1507,2683,63).
\]

Dividing their squares by the corresponding diagonal entries of \(D\)
shows that the largest squared entry of \(\widetilde Q\) is

\nopagebreak[4]

\[
\|\widetilde Q\|_{\max}^2=\frac{63^2}{5272}.
\]

Using the Schur-complement formula from Section 2 gives the following
exact table. Each entry in its second column is
\(\sqrt{5272}\,\|\widetilde S_k\|_{\max}\); the third column identifies
an entry attaining the maximum in the original row and column numbering.

{\def\LTcaptype{none} % do not increment counter
\begin{longtable}[]{@{}lrl@{}}
\toprule\noalign{}
Stage \(k\) & \(\sqrt{5272}\,\|\widetilde S_k\|_{\max}\) & Location \\
\midrule\noalign{}
\endhead
\bottomrule\noalign{}
\endlastfoot
1 & \(63\) & \((1,8)\) \\
2 & \(94\) & \((2,8)\) \\
3 & \(173\) & \((3,8)\) \\
4 & \(338\) & \((4,8)\) \\
5 & \(672\) & \((5,8)\) \\
6 & \(1342\) & \((6,8)\) \\
7 & \(2683\) & \((7,8)\) \\
8 & \(5272\) & \((8,8)\) \\
\end{longtable}
}

These values can be verified directly from the displayed integer
matrices: form \(\widetilde L_{I_k,I_k}T_{I_k,I_k}\), square each entry
in column \(j\), and divide by \(D_{jj}\). In particular the maximum
over all stages occurs at the final pivot and equals \(\sqrt{5272}\).
Hence

\nopagebreak[4]

\[
\rho_{\mathrm{PP}}(\widetilde Q)
=\frac{\sqrt{5272}}{63/\sqrt{5272}}
=\frac{5272}{63}.
\]

\subsection{4. Exact comparison with the specified
candidate}\label{exact-comparison-with-the-specified-candidate}

For completeness, the candidate \(Q_8\) also has an integer-column
description:

\nopagebreak[4]

\[
H_0=
\begin{pmatrix}
1&-5&-8&-12&-16&-16&0&64\\
-1&13&-4&-6&-8&-8&0&32\\
-1&-3&51&-3&-4&-4&0&16\\
-1&-3&-11&169&-2&-2&0&8\\
-1&-3&-11&-43&511&-1&0&4\\
-1&-3&-11&-43&-171&1365&0&2\\
-1&-3&-11&-43&-171&-683&1&1\\
-1&-3&-11&-43&-171&-683&-1&1
\end{pmatrix},
\]

\nopagebreak[4]

\[
D_0=\operatorname{diag}(8,248,3286,36146,349184,2796544,2,5462).
\]

Direct multiplication gives \(H_0^TH_0=D_0\), while \(T_0=L_8^{-1}H_0\)
is upper triangular with positive diagonal

\nopagebreak[4]

\[
(1,8,31,106,341,1024,1,5462).
\]

Thus \(Q_8=H_0D_0^{-1/2}\) has exactly the prescribed QR convention. Its
no-exchange path is again the partial-pivoting path with the first
available row chosen in a tie.

The largest absolute entries of the columns of \(H_0\) are

\nopagebreak[4]

\[
(1,13,51,169,511,1365,1,64),
\]

so exact squared comparisons give

\nopagebreak[4]

\[
\|Q_8\|_{\max}^2=\frac{51^2}{3286},
\]

attained at \((3,3)\). The corresponding active maxima are

{\def\LTcaptype{none} % do not increment counter
\begin{longtable}[]{@{}ll@{}}
\toprule\noalign{}
Stage \(k\) & \(\|S_k^{(0)}\|_{\max}\) \\
\midrule\noalign{}
\endhead
\bottomrule\noalign{}
\endlastfoot
1 & \(51/\sqrt{3286}\) \\
2 & \(96/\sqrt{5462}\) \\
3 & \(176/\sqrt{5462}\) \\
4 & \(344/\sqrt{5462}\) \\
5 & \(684/\sqrt{5462}\) \\
6 & \(1366/\sqrt{5462}\) \\
7 & \(2731/\sqrt{5462}\) \\
8 & \(5462/\sqrt{5462}\) \\
\end{longtable}
}

The table follows either by direct no-exchange elimination of \(H_0\)
with the fixed positive column scalings, or by

\nopagebreak[4]

\[
S_k^{(0)}=(L_8)_{I_k,I_k}(T_0)_{I_k,I_k}(D_0)_{I_k,I_k}^{-1/2}.
\]

Therefore

\nopagebreak[4]

\[
\rho_{\mathrm{PP}}(Q_8)^2
=\frac{5462\cdot3286}{2601}
=\frac{17948132}{2601}.
\]

Finally, the squared growth factors differ by the strictly positive
rational number

\nopagebreak[4]

\[
\left(\frac{5272}{63}\right)^2
-\frac{17948132}{2601}
=\frac{117335164}{1147041}>0.
\]

Both growth factors are positive, so the claimed strict inequality
follows. This proves the theorem. \(\square\)

\subsection{Scope and attribution}\label{scope-and-attribution}

This counterexample refutes the exact extremizer equality asked in
IE-05. It does not determine the true orthogonal supremum in dimension
eight or in general, and it does not refute a separate assertion about
the asymptotic exponential rate or its leading constant.

The extremizer conjecture and the cited element-growth analysis are due
to John Peca-Medlin,
\href{https://arxiv.org/html/2308.16146v2}{\emph{Growth factors of
orthogonal matrices and local behavior of Gaussian elimination with
partial and complete pivoting}, arXiv:2308.16146v2}, Section 3.2 and
Appendix B; published in \href{https://doi.org/10.1137/23M1597733}{SIAM
Journal on Matrix Analysis and Applications 45 (2024), 1599--1620}. The
proof above is self-contained: it does not invoke the paper's extremal
assertions to certify either growth factor.
\end{document}
